% Force exercée par la Terre sur la Station spatiale internationale, à l'échelle 1 cm pour 1,0 × 10^6 N (corrigé du devoir n° 4, exercice 1)
\begin{tikzpicture}[x=1cm,y=1cm]
\shade[ball color=otNavy!35] (0,0) circle (1.6);
\fill[otNavy] (0,0) circle (1.5pt);
\node[pcLegende,anchor=north] at (0,-0.05) {centre O};
\node[pcLegendeG,text=white] at (0,0.75) {Terre};
\draw[otMuted,line width=0.4pt,dash pattern=on 2pt off 2pt] (-16:5.6) arc[start angle=-16,end angle=16,radius=5.6];
\fill[otInk] (5.6,0) circle (2.4pt) node[pcLegende,anchor=south west] {ISS};
\draw[otRed!85!black,line width=1.3pt,-{Stealth[length=2.6mm]}] (5.6,0) -- ++(-3.65,0)
  node[midway,pcLegende,anchor=south,text=otRed!85!black] {$\vec F_{\mathrm{T/ISS}}$ : \qty{3,65e6}{N}};
\node[pcLegende,anchor=north,text=otMuted] at (5.6,-0.15) {échelle des longueurs non respectée};
\end{tikzpicture}
\par\smallskip\legende{La force est dirigée de la station vers le centre de la Terre ; flèche de \qty{3,65}{cm} (\qty{1}{cm} pour \qty{1,0e6}{N}).}
